\documentclass[12pt,a4paper]{article}
\usepackage[top=10mm,left=15mm,right=15mm,bottom=15mm]{geometry}
\usepackage{amsmath,amsthm,amssymb,amsfonts}
\usepackage{mathtools}
\usepackage{booktabs,array}
\usepackage{enumitem}
\usepackage{xcolor}
\usepackage{tocloft}
% Compact TOC
\setlength{\cftbeforesecskip}{2pt}
\setlength{\cftbeforesubsecskip}{0pt}
\renewcommand{\cftsecleader}{\cftdotfill{\cftdotsep}}
\usepackage{hyperref}
\hypersetup{colorlinks=true,linkcolor=blue!70!black,urlcolor=blue!80!black,citecolor=blue!70!black}
\usepackage{cleveref}

% Theorem environments
\theoremstyle{definition}
\newtheorem{theorem}{Theorem}[section]
\newtheorem{lemma}[theorem]{Lemma}
\newtheorem{proposition}[theorem]{Proposition}
\newtheorem{corollary}[theorem]{Corollary}
\newtheorem{conjecture}[theorem]{Conjecture}
\newtheorem{definition}[theorem]{Definition}
\newtheorem{example}[theorem]{Example}
\newtheorem{notation}[theorem]{Notation}
\newtheorem{remark}[theorem]{Remark}

% Non-italic proof label
\makeatletter
\renewenvironment{proof}[1][\proofname]{\par
  \pushQED{\qed}%
  \normalfont \topsep6\p@\@plus6\p@\relax
  \trivlist
  \item[\hskip\labelsep\bfseries #1\@addpunct{.}]\ignorespaces
}{%
  \popQED\endtrivlist\@endpefalse
}
\makeatother

% Custom commands
\newcommand{\R}{\mathbb{R}}
\newcommand{\N}{\mathbb{N}}
\newcommand{\PP}{\mathsf{P}}
\newcommand{\EE}{\mathsf{E}}
\newcommand{\Var}{\operatorname{var}}
\newcommand{\Cov}{\operatorname{cov}}
\newcommand{\tr}{\operatorname{tr}}
\newcommand{\rank}{\operatorname{rank}}
\newcommand{\diag}{\operatorname{diag}}
\newcommand{\dd}{\mathrm{d}}
\newcommand{\ind}{\mathbf{1}}
\newcommand{\caln}{\mathcal{N}}
\newcommand{\erf}{\operatorname{erf}}
\newcommand{\Mn}{M_{n\times m}}

\usepackage{parskip}
\usepackage{etoolbox}
\AtBeginEnvironment{definition}{\vspace{\parskip}}
\AtBeginEnvironment{example}{\vspace{\parskip}}
\AtBeginEnvironment{lemma}{\vspace{\parskip}}
\AtBeginEnvironment{theorem}{\vspace{\parskip}}
\AtBeginEnvironment{corollary}{\vspace{\parskip}}
\AtBeginEnvironment{conjecture}{\vspace{\parskip}}
\AtBeginEnvironment{proposition}{\vspace{\parskip}}
\AtBeginEnvironment{notation}{\vspace{\parskip}}
\AtBeginEnvironment{remark}{\vspace{\parskip}}
\AtBeginEnvironment{exercise}{\vspace{\parskip}}

% Compact section spacing
\usepackage{titlesec}
\titlespacing*{\section}{0pt}{6pt}{3pt}
\titlespacing*{\subsection}{0pt}{4pt}{2pt}
\titlespacing*{\subsubsection}{0pt}{3pt}{1pt}

\begin{document}

% Compact title (no \maketitle top-padding)
\begingroup
\centering
{\LARGE\bfseries The Gaussian Correlation Inequality\par}
\vspace{2pt}
{\normalsize After T.\ Royen (2014) and R.\ Lata{\l}a \& D.\ Matlak (2015)\par}
\vspace{4pt}
\endgroup


%% ====================================================================
\section{Introduction}\label{sec:intro}
%% ====================================================================

\begin{theorem}[Gaussian correlation inequality]\label{thm:gcc_geom}
  For any closed symmetric convex sets $K,L\subseteq\R^d$ and any
  centered Gaussian measure $\mu$ on $\R^d$,
  \begin{equation}\label{eq:gcc}
    \mu(K\cap L)\;\geqslant\;\mu(K)\,\mu(L).
  \end{equation}
\end{theorem}

These notes prove \Cref{thm:gcc_geom} by Royen's analytic argument,
following Lata{\l}a--Matlak's Gaussian-only exposition.  The inputs
are the Laplace transform of a squared Gaussian, a determinant
expansion, dominated convergence, and integration by parts.
Historical background and further remarks are collected in
\Cref{sec:notes}.

\begin{notation}\label{not:conventions}
  \leavevmode
  \begin{itemize}[nosep]
    \item $\caln(0,C)$: centered Gaussian distribution with covariance
      $C$.
    \item $M_{n\times m}$: real $n\times m$ matrices.  $A\succ 0$
      (resp.\ $A\succcurlyeq 0$) means $A$ is symmetric positive
      (semi)definite.
    \item For $A=(a_{ij})_{i,j\leqslant n}$ and
      $J\subseteq[n]:=\{1,\ldots,n\}$,
      $A_J:=(a_{ij})_{i,j\in J}$ is the principal submatrix,
      $|A|$ the determinant, $|J|$ the cardinality.
    \item $I_n$: $n\times n$ identity.
    \item For $\varepsilon\in\{-1,1\}^n$, $x\in(0,\infty)^n$:
      $\sqrt{x}_\varepsilon:=(\varepsilon_i\sqrt{x_i})_{i\leqslant n}$.
  \end{itemize}
\end{notation}

We will use three standard facts without further comment:

\begin{lemma}[Square root]\label{lem:sqrt}
  For any $A\succcurlyeq 0$ there exists a unique $A^{1/2}\succcurlyeq 0$ with
  $(A^{1/2})^2=A$, and $A^{1/2}$ is invertible iff $A\succ 0$.
\end{lemma}

(Spectral theorem: diagonalize $A=O\Lambda O^\top$ with $\Lambda\succcurlyeq 0$,
set $A^{1/2}:=O\Lambda^{1/2}O^\top$; uniqueness is standard.)

\begin{lemma}[Sylvester determinant identity]\label{lem:sylvester}
  For $A_1\in M_{n\times m}$ and $A_2\in M_{m\times n}$,
  $|I_n+A_1A_2|=|I_m+A_2A_1|$.
\end{lemma}

\begin{proof}
  Compute $D:=\bigl|\begin{smallmatrix}I_n & -A_1\\ A_2 &
  I_m\end{smallmatrix}\bigr|$ two ways.  Row-reduce $R_1\leftarrow
  R_1+A_1R_2$ to clear the $(1,1)$-block's contribution to the $(1,2)$
  block: the matrix becomes
  $\bigl(\begin{smallmatrix}I_n+A_1A_2 & 0\\ A_2 & I_m\end{smallmatrix}\bigr)$,
  so $D=|I_n+A_1A_2|$.  Alternatively, row-reduce $R_2\leftarrow
  R_2-A_2R_1$: the matrix becomes
  $\bigl(\begin{smallmatrix}I_n & -A_1\\ 0 & I_m+A_2A_1\end{smallmatrix}\bigr)$,
  so $D=|I_m+A_2A_1|$.
\end{proof}

\begin{lemma}[Strict separation; Hahn--Banach]\label{lem:separation}
  Let $K_1,K_2\subseteq\R^d$ be disjoint closed convex sets with $K_1$
  bounded.  There exist $w\in\R^d\setminus\{0\}$ and $\alpha_1<\alpha_2$
  in $\R$ with $\langle w,x\rangle\leqslant\alpha_1$ for $x\in K_1$ and
  $\langle w,x\rangle\geqslant\alpha_2$ for $x\in K_2$.
\end{lemma}

(Standard corollary of Hahn--Banach in finite dimension; see e.g.\
Rockafellar, \emph{Convex Analysis}, Cor.\ 11.4.2.)


%% ====================================================================
\section{Setup and reduction to rectangles}\label{sec:setup}
%% ====================================================================

\begin{definition}[Centered Gaussian measure]\label{def:gaussian}
  For $C\succ 0$, the measure $\caln(0,C)$ on $\R^n$ has density
  \[
    \phi_C(x)
    \;=\;
    (2\pi)^{-n/2}|C|^{-1/2}
    \exp\bigl(-\tfrac{1}{2}\langle C^{-1}x,x\rangle\bigr).
  \]
\end{definition}

We assume $C\succ 0$ throughout the proof; the extension to singular
$C\succcurlyeq 0$ is deferred to \Cref{note:singular}.

\begin{definition}\label{def:symconvex}
  $K\subseteq\R^n$ is \emph{symmetric} if $K=-K$.  A \emph{symmetric
  slab} is
  $S_{v,t}:=\{x\in\R^n:|\langle v,x\rangle|\leqslant t\}$ for
  $v\in\R^n\setminus\{0\}$, $t>0$.
\end{definition}

\begin{lemma}[Slab approximation]\label{lem:slab_approx}
  Every nonempty \emph{bounded} closed symmetric convex set
  $K\subseteq\R^d$ is a countable intersection of symmetric slabs
  $S_{v,t}$.
\end{lemma}

\begin{proof}
  Fix $R<\infty$ with $K\subseteq B(0,R)$, and let
  $\mathcal V:=\{(v,t)\in\mathbb{Q}^d\times\mathbb{Q}_{>0}:K\subseteq S_{v,t}\}$,
  a countable set; trivially $K\subseteq\bigcap_{(v,t)\in\mathcal V}S_{v,t}$.
  For the reverse inclusion we show that each $y\notin K$ is excluded by
  some slab in $\mathcal V$.

  Fix $y\notin K$ and pick $\varepsilon>0$ with
  $B(y,\varepsilon)\cap K=\varnothing$.  The sets
  $\overline{B(y,\varepsilon/2)}$ (bounded) and $K$ (closed convex) are
  disjoint, so \Cref{lem:separation} yields a separating functional;
  replacing it by its negative if necessary, we obtain
  $w\in\R^d\setminus\{0\}$ and $\alpha_1<\alpha_2$ with
  $\langle w,x\rangle\leqslant\alpha_1$ for $x\in K$ and
  $\langle w,x\rangle\geqslant\alpha_2$ for $x\in\overline{B(y,\varepsilon/2)}$.
  Set $a:=\sup_{x\in K}\langle w,x\rangle\leqslant\alpha_1$.  Since $K$ is
  nonempty, symmetric, and convex, $0=\tfrac12 x+\tfrac12(-x)\in K$ for
  any $x\in K$, so $a\geqslant\langle w,0\rangle=0$; and $K=-K$ gives
  $\sup_{x\in K}|\langle w,x\rangle|\leqslant a$.  The minimum of
  $\langle w,\cdot\rangle$ over $\overline{B(y,\varepsilon/2)}$ is attained at
  $y-\tfrac{\varepsilon}{2}w/|w|$, so
  \[
    \langle w,y\rangle
    =\min_{z\in\overline{B(y,\varepsilon/2)}}\langle w,z\rangle
      +\tfrac{\varepsilon}{2}|w|
    \;\geqslant\;\alpha_2+\tfrac{\varepsilon}{2}|w|
    \;\geqslant\;a+\tfrac{\varepsilon}{2}|w|>0,
  \]
  a strictly positive margin over $a$.

  Now choose $(w',a')\in\mathbb{Q}^d\times\mathbb{Q}_{>0}$ with
  $|w'-w|<\delta$, $|a'-a|<\delta$.  Using $|x|\leqslant R$ on $K$,
  $\sup_{x\in K}|\langle w,x\rangle|\leqslant a$, and
  $\langle w,y\rangle\geqslant a+\tfrac{\varepsilon}{2}|w|$,
  \[
    \sup_{x\in K}|\langle w',x\rangle|
    \leqslant a+\delta R,\qquad
    |\langle w',y\rangle|\geqslant a+\tfrac{\varepsilon}{2}|w|-\delta|y|.
  \]
  Taking $\delta<\dfrac{\varepsilon|w|}{2(|y|+R+1)}$ and
  $a':=a+\delta R+\delta/2\in\mathbb{Q}_{>0}$ then forces $\delta(|y|+R+\tfrac12)
  <\tfrac{\varepsilon}{2}|w|$, which gives both $K\subseteq S_{w',a'}$ and
  $y\notin S_{w',a'}$, as required.
\end{proof}

We now reduce \Cref{thm:gcc_geom} to a statement about slabs.  If $K$ or
$L$ is empty both sides of \eqref{eq:gcc} vanish, so assume both are
nonempty.  We proceed in two steps, each by continuity of the finite measure
$\mu$ (from below for the increasing truncations of the first step, from
above for the decreasing slab-intersections of the second).  First, putting $K_m:=K\cap[-m,m]^d$ and $L_m:=L\cap[-m,m]^d$
(nonempty for $m$ large, bounded closed symmetric convex, increasing to
$K,L$), we get $\mu(K_m)\to\mu(K)$, $\mu(L_m)\to\mu(L)$, and
$\mu(K_m\cap L_m)\to\mu(K\cap L)$, so it suffices to treat \emph{bounded}
$K,L$.  Second, \Cref{lem:slab_approx} writes each bounded set as a
countable intersection of symmetric slabs; intersecting the first $N$ of
them and letting $N\to\infty$, the same continuity reduces the GCI to the
case where $K$ and $L$ are \emph{finite} intersections of symmetric
slabs.  Writing
\[
  K=\bigl\{x : |\langle v_i,x\rangle|\leqslant t_i,\;1\leqslant i\leqslant n_1\bigr\},
  \qquad
  L=\bigl\{x : |\langle v_i,x\rangle|\leqslant t_i,\;n_1{+}1\leqslant i\leqslant n\bigr\},
\]
setting $n=n_1+n_2$ and $X_i:=\langle v_i,G\rangle$ for $G\sim\mu$,
the GCI reduces to:

\begin{theorem}[GCI, analytic form]\label{thm:gcc2}
  Let $X=(X_1,\ldots,X_n)$ be centered Gaussian with $C\succ 0$,
  $n=n_1+n_2$.  For any $t_1,\ldots,t_n>0$,
  \begin{equation}\label{eq:gcc2}
    \PP(|X_i|\leqslant t_i,\;i\leqslant n)
    \;\geqslant\;
    \PP(|X_i|\leqslant t_i,\;i\leqslant n_1)\,
    \PP(|X_i|\leqslant t_i,\;n_1{<}i\leqslant n).
  \end{equation}
\end{theorem}


%% ====================================================================
\section{Auxiliary lemmas}\label{sec:lemmas}
%% ====================================================================

\subsection{Laplace transform of squared Gaussian vectors}

\begin{lemma}\label{lem:Lap}
  For $X\sim\caln(0,C)$ in $\R^n$ and $\lambda_i\geqslant 0$,
  \[
    \EE\exp\Bigl(-\sum_{i=1}^n\lambda_iX_i^2\Bigr)
    \;=\;|I_n+2\Lambda C|^{-1/2},
    \qquad\Lambda:=\diag(\lambda_1,\ldots,\lambda_n).
  \]
\end{lemma}

\begin{proof}
  \leavevmode
  \begin{enumerate}
    \item Standard normal quadratic.  For $Y\sim\caln(0,I_n)$ and a
      symmetric matrix $B$ with $I_n-2B\succ 0$, combine the density of
      $Y$ with the factor $e^{\langle BY,Y\rangle}$: the exponent is
      $-\tfrac12\langle y,y\rangle+\langle By,y\rangle
      =-\tfrac12\langle(I_n-2B)y,y\rangle$, so
      \begin{equation}\label{eq:gauss_quad}
        \EE e^{\langle BY,Y\rangle}
        =(2\pi)^{-n/2}\int_{\R^n}e^{-\frac12\langle(I_n-2B)y,y\rangle}\,\dd y
        =|I_n-2B|^{-1/2},
      \end{equation}
      using $\int_{\R^n}e^{-\frac12\langle Ay,y\rangle}\dd y=(2\pi)^{n/2}|A|^{-1/2}$
      (change variables $y=A^{-1/2}z$; $A=I_n-2B$).

    \item Reduction to the standard case.  Set $A:=C^{1/2}$
      (\Cref{lem:sqrt}); then $C=AA^\top$ and $X:=AY\sim\caln(0,C)$.
      Taking $B:=-A^\top\Lambda A$ in \eqref{eq:gauss_quad}: since
      $\lambda_i\geqslant 0$ gives $\Lambda\succcurlyeq 0$, hence $B\preccurlyeq 0$
      and $I_n-2B\succcurlyeq I_n\succ 0$.  The exponent becomes
      $\langle BY,Y\rangle=-\langle\Lambda AY,AY\rangle
      =-\sum_i\lambda_i(AY)_i^2=-\sum_i\lambda_iX_i^2$, so
      \[
        \EE\exp\Bigl(-\sum_i\lambda_iX_i^2\Bigr)
        =|I_n+2A^\top\Lambda A|^{-1/2}.
      \]

    \item Commuting the factors.  Apply \Cref{lem:sylvester} with
      $A_1:=A^\top$, $A_2:=2\Lambda A$:
      $|I_n+2A^\top\Lambda A|=|I_n+2\Lambda AA^\top|=|I_n+2\Lambda C|$.
  \end{enumerate}
\end{proof}


\subsection{A determinant identity}

\begin{lemma}\label{lem:det}
  \leavevmode
  \begin{enumerate}[label=(\roman*),nosep]
    \item For $A\in M_{n\times n}$,
    \begin{equation}\label{eq:det1}
      |I_n+A|\;=\;1+\sum_{\varnothing\ne J\subseteq[n]}|A_J|.
    \end{equation}
    \item For $A=\bigl(\begin{smallmatrix}A_{11}&A_{12}\\A_{21}&A_{22}
      \end{smallmatrix}\bigr)$ with $A_{11},A_{22}\succ 0$,
    \begin{equation}\label{eq:det2}
      |A|=|A_{11}||A_{22}|\cdot
      \bigl|I_{n_1}-A_{11}^{-1/2}A_{12}A_{22}^{-1}A_{21}A_{11}^{-1/2}\bigr|.
    \end{equation}
    If moreover $A\succ 0$, then
    $0\preccurlyeq A_{11}^{-1/2}A_{12}A_{22}^{-1}A_{21}
    A_{11}^{-1/2}\preccurlyeq I_{n_1}$.
  \end{enumerate}
\end{lemma}

\begin{proof}
  \leavevmode
  \begin{enumerate}[label=(\roman*)]
    \item By Leibniz,
      $|I_n+A|=\sum_\sigma\mathrm{sgn}(\sigma)\prod_i
      (\delta_{i\sigma(i)}+a_{i\sigma(i)})$.  Expanding the product, each
      summand corresponds to a choice of $J\subseteq[n]$ (positions where
      $a_{i\sigma(i)}$ is selected); non-zero contributions require
      $\sigma(i)=i$ for $i\notin J$.  The $\sigma$'s fixing $[n]\setminus J$
      correspond to permutations of $J$, giving $\sum_J|A_J|$; the $J=\varnothing$
      term is $1$.

    \item Let $A_{11}^{1/2},A_{22}^{1/2}$ be the positive
      symmetric square roots (\Cref{lem:sqrt}), invertible since
      $A_{11},A_{22}\succ 0$.
      Setting $P:=A_{11}^{-1/2}A_{12}A_{22}^{-1/2}$ and
      $Q:=A_{22}^{-1/2}A_{21}A_{11}^{-1/2}$, a direct block multiplication
      gives
      \[
        A=\diag(A_{11}^{1/2},A_{22}^{1/2})\,
        \bigl(\begin{smallmatrix}I&P\\Q&I\end{smallmatrix}\bigr)\,
        \diag(A_{11}^{1/2},A_{22}^{1/2}),
      \]
      so $|A|=|A_{11}||A_{22}|\cdot\bigl|
      \begin{smallmatrix}I&P\\Q&I\end{smallmatrix}\bigr|$.
      Row-reduce by $R_1\leftarrow R_1-PR_2$ (determinant invariant):
      \[
        \bigl(\begin{smallmatrix}I&P\\Q&I\end{smallmatrix}\bigr)
        \;\to\;
        \bigl(\begin{smallmatrix}I-PQ&0\\Q&I\end{smallmatrix}\bigr),
      \]
      block-triangular with determinant $|I-PQ|$.  Substituting
      $PQ=A_{11}^{-1/2}A_{12}A_{22}^{-1}A_{21}A_{11}^{-1/2}$ gives
      \eqref{eq:det2}.

      When $A\succ 0$: for all $s,t\in\R$ and $x\in\R^{n_1}$,
      $y\in\R^{n_2}$,
      \[
        \langle A(sx,ty)^\top,(sx,ty)^\top\rangle
        =s^2\langle A_{11}x,x\rangle+2st\langle A_{21}x,y\rangle
        +t^2\langle A_{22}y,y\rangle\geqslant 0.
      \]
      Nonnegativity of this quadratic in $(s,t)$ forces
      $\langle A_{21}x,y\rangle^2\leqslant\langle A_{11}x,x\rangle
      \langle A_{22}y,y\rangle$.
      Setting $x:=A_{11}^{-1/2}u$, $y:=A_{22}^{-1/2}v$:
      $\langle Qu,v\rangle^2\leqslant|u|^2|v|^2$, i.e.,
      $\|Q\|_{\mathrm{op}}\leqslant 1$, equivalently
      $0\preccurlyeq Q^\top Q\preccurlyeq I_{n_1}$.  For $A$ symmetric, $Q=P^\top$, so
      $A_{11}^{-1/2}A_{12}A_{22}^{-1}A_{21}A_{11}^{-1/2}=PP^\top=Q^\top Q$.
  \end{enumerate}
\end{proof}


\subsection{Density of $Z(\tau)$ and the interchange lemma}

Partition
\begin{equation}\label{eq:partition}
  C=\begin{pmatrix}C_{11}&C_{12}\\C_{21}&C_{22}\end{pmatrix},
  \quad C_{ij}\in M_{n_i\times n_j},\;n_1+n_2=n,
\end{equation}
and define the interpolating family
\begin{equation}\label{eq:Ctau}
  C(\tau):=\begin{pmatrix}C_{11}&\tau C_{12}\\
    \tau C_{21}&C_{22}\end{pmatrix}
  =\tau C(1)+(1-\tau)C(0),\qquad\tau\in[0,1].
\end{equation}
Since $C(0)=\diag(C_{11},C_{22})\succ 0$ and $C(1)=C\succ 0$, each
$C(\tau)\succ 0$ by convexity.  Let $X(\tau)\sim\caln(0,C(\tau))$ and
$Z_i(\tau):=\tfrac{1}{2}X_i(\tau)^2$.

\begin{lemma}[Density and interchange]\label{lem:diff}
  Let $f(x,\tau)$ denote the density of $Z(\tau)$ on $(0,\infty)^n$.
  \begin{enumerate}[label=(\roman*),nosep]
    \item
    \begin{equation}\label{eq:fxtau}
      f(x,\tau)
      =|C(\tau)|^{-1/2}(4\pi)^{-n/2}
      \frac{1}{\sqrt{x_1\cdots x_n}}
      \sum_{\varepsilon\in\{-1,1\}^n}
      e^{-\langle C(\tau)^{-1}\sqrt{x}_\varepsilon,
        \sqrt{x}_\varepsilon\rangle}
      \ind_{(0,\infty)^n}(x).
    \end{equation}
    \item For any Borel $K\subseteq[0,\infty)^n$ and
    $\lambda_i\geqslant 0$,
    \begin{equation}\label{eq:interchange}
      \int_K e^{-\sum_i\lambda_ix_i}\,
      \frac{\partial}{\partial\tau}f(x,\tau)\,\dd x
      \;=\;
      \frac{\partial}{\partial\tau}\int_K e^{-\sum_i\lambda_ix_i}\,
      f(x,\tau)\,\dd x.
    \end{equation}
  \end{enumerate}
\end{lemma}

\begin{proof}
  \leavevmode
  \begin{enumerate}[label=(\roman*)]
    \item The density of $X(\tau)$ is $\phi_{C(\tau)}$ from
      \Cref{def:gaussian}.  On each sign orthant
      $\R^n_\varepsilon:=\{y:\text{sgn}(y_i)=\varepsilon_i\}$, the map
      $y\mapsto x_i:=\tfrac12 y_i^2$ is a diffeomorphism onto
      $(0,\infty)^n$ with inverse $y_i=\varepsilon_i\sqrt{2x_i}$ and
      Jacobian $\dd y_i/\dd x_i=\varepsilon_i/\sqrt{2x_i}$; absolute
      Jacobian $\prod_i(2x_i)^{-1/2}$.  Summing over the $2^n$ sign
      patterns and using
      $\langle C(\tau)^{-1}y,y\rangle
      =2\langle C(\tau)^{-1}\sqrt{x}_\varepsilon,\sqrt{x}_\varepsilon\rangle$
      (by $y_i=\varepsilon_i\sqrt{2x_i}$):
      \begin{align*}
        f(x,\tau)
        &=\sum_\varepsilon\phi_{C(\tau)}(\sqrt{2x}_\varepsilon)
          \cdot\prod_i(2x_i)^{-1/2}\\
        &=|C(\tau)|^{-1/2}(2\pi)^{-n/2}2^{-n/2}(x_1\cdots x_n)^{-1/2}
          \sum_\varepsilon e^{-\langle C(\tau)^{-1}\sqrt{x}_\varepsilon,
          \sqrt{x}_\varepsilon\rangle},
      \end{align*}
      which is \eqref{eq:fxtau} after combining $(2\pi)^{-n/2}2^{-n/2}=(4\pi)^{-n/2}$.

    \item We build a $\tau$-uniform integrable dominating function for
      $\partial_\tau f$ via four steps.
      \begin{enumerate}
        \item Eigenvalue lower bound.  Each $C(\tau)\succ 0$, and
          $\tau\mapsto\lambda_{\min}(C(\tau)^{-1})=1/\lambda_{\max}(C(\tau))$
          is continuous and positive on compact $[0,1]$, so bounded
          below by some $a>0$.  Then
          $\langle C(\tau)^{-1}\sqrt{x}_\varepsilon,\sqrt{x}_\varepsilon\rangle
          \geqslant a|\sqrt{x}_\varepsilon|^2=a\sum_ix_i$.

        \item Derivative formula.  Differentiating
          $C(\tau)C(\tau)^{-1}=I_n$,
          $\tfrac{\dd}{\dd\tau}C(\tau)^{-1}=-C(\tau)^{-1}C'(\tau)C(\tau)^{-1}$,
          with $C'(\tau)=C(1)-C(0)$ constant.  Since
          $\tau\mapsto C(\tau)^{-1}$ is continuous on compact $[0,1]$,
          both $|C(\tau)|^{-1/2}$ and
          $\|\tfrac{\dd}{\dd\tau}C(\tau)^{-1}\|_{\mathrm{op}}$ are
          bounded, say by $M,b<\infty$.  Hence
          $|\partial_\tau|C(\tau)|^{-1/2}|\leqslant M$ and
          $|\langle\tfrac{\dd}{\dd\tau}C(\tau)^{-1}\sqrt{x}_\varepsilon,
          \sqrt{x}_\varepsilon\rangle|\leqslant b|\sqrt{x}_\varepsilon|^2
          =b\sum_ix_i$.

        \item Dominating function.  Differentiating \eqref{eq:fxtau}
          in $\tau$ (product rule on the prefactor $|C(\tau)|^{-1/2}$
          and the exponent), each of the $2^n$ sign-orthant summands
          is bounded by
          \[
            \bigl(M+b\textstyle\sum_ix_i\bigr)
            (4\pi)^{-n/2}(x_1\cdots x_n)^{-1/2}e^{-a\sum_ix_i}.
          \]
          Summing over $\varepsilon$ and absorbing constants,
          \[
            \sup_{\tau\in[0,1]}|\partial_\tau f(x,\tau)|
            \leqslant g(x):=c\,\frac{1+\sum_ix_i}{\sqrt{x_1\cdots x_n}}
            e^{-a\sum_ix_i}
          \]
          for some $c=c(C,n)$.  Since $1+\sum_ix_i\leqslant\prod_i(1+x_i)$,
          we have $g(x)\leqslant c\prod_i x_i^{-1/2}(1+x_i)e^{-ax_i}$,
          which factorises; as $\int_0^\infty x^{-1/2}(1+x)e^{-ax}\,\dd x<\infty$,
          Tonelli gives $\int g<\infty$.

        \item Interchange.  For any Borel $K\subseteq[0,\infty)^n$ and
          $\lambda_i\geqslant 0$,
          $e^{-\sum_i\lambda_ix_i}|\partial_\tau f|\leqslant g$
          uniformly in $\tau$.  Dominated convergence applied to the
          difference quotient
          $(f(x,\tau+h)-f(x,\tau))/h\to\partial_\tau f(x,\tau)$ gives
          \eqref{eq:interchange}.
      \end{enumerate}
  \end{enumerate}
\end{proof}


\subsection{The noncentral gamma and the $h$-family}

\begin{definition}\label{def:ncgamma}
  For $\alpha>0$, $y\geqslant 0$, the \emph{noncentral gamma density} is
  \[
    g_\alpha(x,y)
    :=e^{-x-y}\sum_{k=0}^\infty
    \frac{x^{k+\alpha-1}}{\Gamma(k+\alpha)}\,\frac{y^k}{k!},
    \qquad x>0.
  \]
\end{definition}

\begin{remark}\label{rem:ncgamma_poisson}
  Equivalently,
  \[
    g_\alpha(x,y)=\sum_{k\geqslant 0}e^{-y}\tfrac{y^k}{k!}g_{\alpha+k}(x),
  \]
  a Poisson$(y)$ mixture of central gamma$(\alpha+k)$ densities; so
  $g_\alpha(\cdot,y)$ is a probability density for each $y\geqslant 0$.
\end{remark}

\begin{definition}\label{def:h_family}
  For a scale parameter $\mu>0$ (unrelated to the Gaussian measure of
  \Cref{thm:gcc_geom}), $\alpha\in(0,\infty)^n$, and
  $Y=(Y_1,\ldots,Y_n)$ with $Y_i\geqslant 0$ a.s.,
  \[
    h_{\alpha,\mu,Y}(x)
    :=\EE\prod_{i=1}^n
      \tfrac{1}{\mu}\,g_{\alpha_i}\!\bigl(\tfrac{x_i}{\mu},Y_i\bigr),
    \qquad x\in(0,\infty)^n.
  \]
\end{definition}

\begin{lemma}[Properties of the $h$-family]\label{lem:proph}
  Write $h_\alpha:=h_{\alpha,\mu,Y}$.
  \begin{enumerate}[label=(\roman*)]
    \item $h_\alpha\geqslant 0$ and $\int h_\alpha\,\dd x=1$.
    \item If $\alpha_i>1$, then $h_\alpha(x)\to 0$ as $x_i\to 0+$,
      and
      \begin{equation}\label{eq:deriv_h}
        \frac{\partial}{\partial x_i}h_\alpha=\tfrac{1}{\mu}\bigl(h_{\alpha-e_i}-h_\alpha\bigr).
      \end{equation}
    \item If $\alpha\in(1,\infty)^n$, then for any $J\subseteq[n]$
      and $\lambda_i\geqslant 0$,
      $\partial^{|J|}h_\alpha/\partial x_J\in L_1((0,\infty)^n)$ and
      \begin{equation}\label{eq:proph_iii}
        \int e^{-\sum_i\lambda_ix_i}
          \frac{\partial^{|J|}}{\partial x_J}h_\alpha\,\dd x
        \;=\;
        \Bigl(\prod_{i\in J}\lambda_i\Bigr)
        \int e^{-\sum_i\lambda_ix_i}h_\alpha\,\dd x.
      \end{equation}
  \end{enumerate}
\end{lemma}

\begin{proof}
  \leavevmode
  \begin{enumerate}[label=(\roman*)]
    \item $\int_0^\infty\mu^{-1}g_\alpha(x_i/\mu,Y_i)\,\dd x_i
      =\int_0^\infty g_\alpha(u,Y_i)\,\dd u=1$ for each $\omega$
      (\Cref{rem:ncgamma_poisson}); Fubini gives
      $\int h_\alpha=\EE[1]=1$.

    \item We split into two parts.
      \begin{enumerate}
        \item Gamma bound.  $\log\Gamma$ is strictly convex on
          $(0,\infty)$, and since $\Gamma(1)=\Gamma(2)=1$, $\Gamma$ has
          a unique minimum at some $s_0\in(1,2)$ with $\Gamma(s_0)>\tfrac12$
          (numerically $s_0\approx 1.462$, $\Gamma(s_0)\approx 0.886$).
          Hence
          \begin{equation}\label{eq:gamma_lb}
            \Gamma(s)\geqslant\tfrac12\quad\text{for all }s>0,
          \end{equation}
          giving $\Gamma(\alpha)^{-1}\leqslant 2$ for the $k=0$ term.
          For $k\geqslant 1$, iterating $\Gamma(s+1)=s\Gamma(s)$ gives
          \[
            \Gamma(k+\alpha)
            =\Gamma(1+\alpha)\prod_{j=1}^{k-1}(j+\alpha)
            \;\geqslant\;\tfrac12\prod_{j=1}^{k-1}j
            =\tfrac12(k-1)!,
          \]
          where we used $\Gamma(1+\alpha)\geqslant\tfrac12$ and
          $j+\alpha\geqslant j$ (the product is empty for $k=1$).
          Substituting into
          $g_\alpha(x,y)$:
          \begin{align*}
            g_\alpha(x,y)
            &\leqslant e^{-x-y}\Bigl[2x^{\alpha-1}
              +\sum_{k\geqslant 1}\tfrac{2x^{k+\alpha-1}y^k}{(k-1)!k!}\Bigr]\\
            &=2x^{\alpha-1}e^{-x-y}
              +2x^\alpha e^{-x}\underbrace{\sum_{k\geqslant 1}\tfrac{x^{k-1}}{(k-1)!}
              \cdot\tfrac{y^ke^{-y}}{k!}}_{\leqslant\;e^x\cdot 1}
            \leqslant 2x^{\alpha-1}(e^{-x}+x),
          \end{align*}
          using $y^ke^{-y}/k!\leqslant\sup_y(y^ke^{-y}/k!)\leqslant 1$
          (term in a Poisson distribution), and
          $\sum_{k\geqslant 1}x^{k-1}/(k-1)!=e^x$.  Thus
          \begin{equation}\label{eq:gbound}
            g_\alpha(x,y)\leqslant 2x^{\alpha-1}(e^{-x}+x),
          \end{equation}
          and substituting into \Cref{def:h_family} (note $Y_i$ drops out),
          \begin{equation}\label{eq:hbound}
            h_\alpha(x)\leqslant(2/\mu)^n\prod_{i}(x_i/\mu)^{\alpha_i-1}(1+x_i/\mu).
          \end{equation}
          When $\alpha_i>1$, the factor $(x_i/\mu)^{\alpha_i-1}\to 0$
          as $x_i\to 0+$, so $h_\alpha(x)\to 0$.

        \item Derivative formula.  Differentiating the series in
          \Cref{def:ncgamma} term by term (using
          $\Gamma(k+\alpha)=(k+\alpha-1)\Gamma(k+\alpha-1)$; for the
          $k=0$ term this reads $\Gamma(\alpha)=(\alpha-1)\Gamma(\alpha-1)$,
          which is exactly where $\alpha>1$ is needed),
          \begin{equation}\label{eq:g_deriv}
            \partial_x g_\alpha(x,y)=g_{\alpha-1}(x,y)-g_\alpha(x,y),
            \qquad\alpha>1,
          \end{equation}
          with convergence locally uniform by \eqref{eq:gbound} applied
          to $g_{\alpha-1}$ (valid since $\alpha-1>0$).  Inserting into
          \Cref{def:h_family}: the integrand and its $x_i$-derivative are
          bounded by an integrable function of $(x,\omega)$ uniformly in
          $x_i$ on compact subsets of $(0,\infty)$ (use \eqref{eq:gbound}
          for $g_{\alpha_i-1}(\cdot,Y_i)$ and $g_{\alpha_i}(\cdot,Y_i)$),
          justifying dominated convergence to differentiate under $\EE$,
          which gives \eqref{eq:deriv_h}.
      \end{enumerate}

    \item Iterating \eqref{eq:deriv_h},
      \begin{equation}\label{eq:derivJ_h}
        \frac{\partial^{|J|}}{\partial x_J}h_\alpha
        =\frac{1}{\mu^{|J|}}\sum_{\delta\in\{0,1\}^J}(-1)^{|J|-|\delta|}
          h_{\alpha-\sum_{i\in J}\delta_ie_i},
      \end{equation}
      each summand in $L_1$ by (i) (all $\alpha_i-\delta_i\geqslant\alpha_i-1>0$).
      Integrate by parts in each $x_j$, $j\in J$, in turn.  At the step
      removing $\partial_{x_j}$, the function being differentiated is a
      finite combination of densities $h_{\alpha'}$ (via
      \eqref{eq:derivJ_h}) in which coordinate $j$ is \emph{not} among
      the lowered ones, so its $j$-th shape parameter is still
      $\alpha_j>1$; by (i) each such $h_{\alpha'}$, and by
      \eqref{eq:derivJ_h} its $x_j$-partial, lies in $L_1$.  Hence, for
      a.e.\ fixing of the remaining coordinates, $x_j\mapsto h_{\alpha'}$
      is absolutely continuous with itself and its derivative integrable;
      such a function tends to $0$ as $x_j\to\infty$ (an $L_1$ function on
      $(0,\infty)$ with integrable derivative has limit $0$ at infinity),
      and to $0$ as $x_j\to 0+$ by (ii) (since $\alpha'_j=\alpha_j>1$).
      Both boundary contributions therefore vanish (the bounded factor
      $e^{-\sum_i\lambda_ix_i}\in(0,1]$ does not affect this), and each
      integration by parts brings down a factor $\lambda_j$, giving
      \eqref{eq:proph_iii}.
  \end{enumerate}
\end{proof}


\subsection{The family $h_{k,C}$ and its Laplace transform}
  \label{sec:hkC}

\begin{definition}\label{def:hkC}
  For $C\succ 0$ and $\mu>0$ with $C-\mu I_n\succ 0$, set
  $A:=(C-\mu I_n)^{1/2}$ (\Cref{lem:sqrt}, which is well-defined and
  invertible), so that $C=\mu I_n+AA^\top$.  Let $(g_j^{(l)})$,
  $j\leqslant n$, $l\leqslant k$, be i.i.d.\ $\caln(0,1)$, and
  \begin{equation}\label{eq:defhkC}
    Y_i:=\sum_{l=1}^k\Bigl(\sum_{j=1}^n
      \tfrac{1}{\sqrt{2\mu}}g_j^{(l)}a_{ij}\Bigr)^2,
    \qquad
    h_{k,C}:=h_{(k/2,\ldots,k/2),\mu,Y}.
  \end{equation}
\end{definition}

\begin{lemma}\label{lem:LaplhkC}
  For $\lambda_i\geqslant 0$,
  $\int e^{-\sum_i\lambda_ix_i}h_{k,C}\,\dd x
  =|I_n+\Lambda C|^{-k/2}$.
\end{lemma}

\begin{proof}
  \leavevmode
  \begin{enumerate}
    \item Scalar Laplace transform.  Substitute $u:=x/\mu$ and expand
      the series in \Cref{def:ncgamma}:
      \begin{align*}
        \int_0^\infty\mu^{-1}e^{-\lambda x}g_\alpha(x/\mu,y)\,\dd x
        &=\int_0^\infty e^{-\lambda\mu u-u-y}
          \sum_{k\geqslant 0}\tfrac{u^{k+\alpha-1}}{\Gamma(k+\alpha)}
          \tfrac{y^k}{k!}\,\dd u\\
        &=e^{-y}\sum_{k\geqslant 0}\tfrac{y^k}{k!}(1+\mu\lambda)^{-(k+\alpha)}\\
        &=(1+\mu\lambda)^{-\alpha}\exp\bigl(-\tfrac{\mu\lambda}{1+\mu\lambda}y\bigr),
      \end{align*}
      using $\int_0^\infty u^{k+\alpha-1}e^{-(1+\mu\lambda)u}\,\dd u
      =\Gamma(k+\alpha)(1+\mu\lambda)^{-(k+\alpha)}$ (Fubini-term-by-term
      justified by absolute convergence).  So
      \begin{equation}\label{eq:1d_lt}
        \int_0^\infty\mu^{-1}e^{-\lambda x}g_\alpha(x/\mu,y)\,\dd x
        =(1+\mu\lambda)^{-\alpha}\exp\bigl(-\tfrac{\mu\lambda}{1+\mu\lambda}y\bigr).
      \end{equation}

    \item Product form.  Applying \eqref{eq:1d_lt} coordinatewise with
      $\alpha_i=k/2$ and Fubini (nonneg integrand),
      \begin{equation}\label{eq:hkC_step1}
        \int e^{-\sum_i\lambda_ix_i}h_{k,C}\,\dd x
        =\prod_i(1+\mu\lambda_i)^{-k/2}\cdot
        \EE\exp\Bigl(-\sum_i\tfrac{\mu\lambda_i}{1+\mu\lambda_i}Y_i\Bigr).
      \end{equation}

    \item Evaluate the expectation.  From \eqref{eq:defhkC}, set
      $X^{(l)}:=\tfrac{1}{\sqrt{2\mu}}Ag^{(l)}$, so
      $X^{(l)}\sim\caln(0,\tfrac{1}{2\mu}AA^\top)$ i.i.d.\ across $l$
      and $Y_i=\sum_{l=1}^k(X_i^{(l)})^2$.  By independence across $l$
      and \Cref{lem:Lap} applied to each $X^{(l)}$ with
      $\tilde\Lambda:=\diag(\tfrac{\mu\lambda_i}{1+\mu\lambda_i})$,
      \[
        \EE\exp\Bigl(-\sum_i\tfrac{\mu\lambda_i}{1+\mu\lambda_i}Y_i\Bigr)
        =\Bigl(\EE\exp(-\langle\tilde\Lambda X^{(1)},X^{(1)}\rangle)\Bigr)^k
        =|I_n+\tfrac{1}{\mu}\tilde\Lambda AA^\top|^{-k/2}.
      \]

    \item Combine.  Inserting step~3 into \eqref{eq:hkC_step1} and
      using $\prod_i(1+\mu\lambda_i)^{-k/2}=|I_n+\mu\Lambda|^{-k/2}$
      (since $I_n+\mu\Lambda$ is diagonal),
      \[
        \int e^{-\sum_i\lambda_ix_i}h_{k,C}\,\dd x
        =\bigl(|I_n+\mu\Lambda|\cdot
          |I_n+\tfrac{1}{\mu}\tilde\Lambda AA^\top|\bigr)^{-k/2}.
      \]
      The diagonal identity
      $\tilde\Lambda_i=\mu\lambda_i/(1+\mu\lambda_i)$ rearranges to
      $(I_n+\mu\Lambda)\tilde\Lambda=\mu\Lambda$.  Hence, using the
      multiplicativity of $\det$,
      \begin{align*}
        |I_n+\mu\Lambda|\cdot|I_n+\tfrac{1}{\mu}\tilde\Lambda AA^\top|
        &=|(I_n+\mu\Lambda)(I_n+\tfrac{1}{\mu}\tilde\Lambda AA^\top)|\\
        &=|I_n+\mu\Lambda
          +\tfrac{1}{\mu}(I_n+\mu\Lambda)\tilde\Lambda\,AA^\top|\\
        &=|I_n+\mu\Lambda+\Lambda AA^\top|
        =|I_n+\Lambda(\mu I_n+AA^\top)|=|I_n+\Lambda C|,
      \end{align*}
      giving the claim.
  \end{enumerate}
\end{proof}


\subsection{Uniqueness of the Laplace transform}

\begin{theorem}\label{thm:laplace_inj}
  If $\varphi_1,\varphi_2\in L_1((0,\infty)^n)$ satisfy
  $\int e^{-\sum_i\lambda_ix_i}\varphi_1\,\dd x
  =\int e^{-\sum_i\lambda_ix_i}\varphi_2\,\dd x$ for all
  $\lambda\in[0,\infty)^n$, then $\varphi_1=\varphi_2$ a.e.
\end{theorem}

\begin{proof}
  Set $\varphi:=\varphi_1-\varphi_2\in L_1((0,\infty)^n)$; the
  hypothesis becomes $\int e^{-\sum_i\lambda_ix_i}\varphi\,\dd x=0$
  for all $\lambda\in[0,\infty)^n$.
  \begin{enumerate}
    \item Change of variables.  With $y_i:=e^{-x_i}$, the map
      $(0,\infty)^n\to(0,1]^n$ has Jacobian $\prod_iy_i^{-1}$, so
      setting
      $\tilde\varphi(y):=\varphi(-\log y)\prod_iy_i^{-1}\in L_1((0,1]^n)$
      and noting $e^{-\lambda_ix_i}=y_i^{\lambda_i}$,
      \begin{equation}\label{eq:moment_zero}
        \int_{(0,1]^n}y^\lambda\tilde\varphi(y)\,\dd y=0
        \qquad\text{for all }\lambda\in[0,\infty)^n.
      \end{equation}

    \item From polynomials to continuous functions.  Restricting
      \eqref{eq:moment_zero} to $\lambda\in\mathbb{Z}_{\geqslant 0}^n$
      and taking finite linear combinations,
      $\int_{[0,1]^n}P(y)\tilde\varphi(y)\,\dd y=0$ for every polynomial
      $P$ (extend $\tilde\varphi$ by zero at $y_i=0$).  By
      Stone--Weierstrass, polynomials are uniformly dense in
      $C([0,1]^n)$; given $\psi\in C([0,1]^n)$ and $\varepsilon>0$,
      pick $P$ with $\|\psi-P\|_\infty<\varepsilon$, so
      \[
        \Bigl|\int\psi\tilde\varphi\Bigr|
        =\Bigl|\int(\psi-P)\tilde\varphi\Bigr|
        \leqslant\varepsilon\|\tilde\varphi\|_1.
      \]
      Sending $\varepsilon\to 0$:
      $\int_{[0,1]^n}\psi\tilde\varphi\,\dd y=0$ for every
      $\psi\in C([0,1]^n)$.

    \item From continuous functions to $\tilde\varphi=0$ a.e.  Let
      $E:=\{\tilde\varphi>0\}$.  By inner/outer regularity of Lebesgue
      measure, there exist compact $K_m\subseteq E$ and open
      $U_m\supseteq E$ with $|U_m\setminus K_m|\to 0$; Urysohn gives
      $\psi_m\in C([0,1]^n)$ with
      $\mathbf{1}_{K_m}\leqslant\psi_m\leqslant\mathbf{1}_{U_m}$.  Then
      $\psi_m\to\mathbf{1}_E$ a.e.\ along a subsequence, and
      $|\psi_m\tilde\varphi|\leqslant|\tilde\varphi|\in L_1$, so
      dominated convergence gives
      $\int_E\tilde\varphi=\lim\int\psi_m\tilde\varphi=0$.  Similarly
      for $\{\tilde\varphi<0\}$, so $\int|\tilde\varphi|=0$ and
      $\tilde\varphi=0$ a.e., hence $\varphi=0$ a.e.
  \end{enumerate}
\end{proof}


%% ====================================================================
\section{Proof of \Cref{thm:gcc2}}\label{sec:proof}
%% ====================================================================

\subsection{Reduction to monotonicity}

With $s_i:=\tfrac{1}{2}t_i^2$, $K:=[0,s_1]\times\cdots\times[0,s_n]$,
and $Z(\tau)$ as above, \eqref{eq:gcc2} is equivalent to
$\PP(Z(1)\in K)\geqslant\PP(Z(0)\in K)$.  At $\tau=0$,
$C(0)=\diag(C_{11},C_{22})$ is block-diagonal, so the two blocks are
independent and the RHS factorizes.  At $\tau=1$ we recover $C$.
Hence \Cref{thm:gcc2} follows from:

\begin{proposition}\label{prop:key}
  $\tau\mapsto\int_K f(x,\tau)\,\dd x$ is non-decreasing on $[0,1]$.
\end{proposition}

\subsection{Laplace-transform computation of $\partial_\tau f$}

\Cref{prop:key} amounts to $\int_K\partial_\tau f(x,\tau)\,\dd x\geqslant 0$
for each $\tau\in(0,1)$ (\Cref{lem:diff}(ii) with $\lambda=0$ justifies
$\partial_\tau\int_K f=\int_K\partial_\tau f$).  We obtain this by:
\begin{enumerate}[label=(\alph*)]
  \item computing the Laplace transform of $\partial_\tau f$ and
    matching it with that of a sum of derivatives of $h_\tau$;
  \item using Laplace injectivity (\Cref{thm:laplace_inj}) to get an
    a.e.\ pointwise identity;
  \item integrating this identity over $K$.
\end{enumerate}

\paragraph{Step (a).}  Since $Z_i(\tau)=\tfrac12 X_i(\tau)^2$, applying
\Cref{lem:Lap} to $X(\tau)\sim\caln(0,C(\tau))$ with $\lambda_i/2$ in
place of $\lambda_i$,
\[
  \int_{[0,\infty)^n}e^{-\sum_i\lambda_ix_i}f(x,\tau)\,\dd x
  =\EE e^{-\sum_i\lambda_iZ_i(\tau)}
  =|I_n+\Lambda C(\tau)|^{-1/2}.
\]
Differentiating in $\tau$, with \Cref{lem:diff}(ii) (taking $K=[0,\infty)^n$)
justifying the interchange:
\begin{equation}\label{eq:lt_deriv}
  \int_{[0,\infty)^n}e^{-\sum_i\lambda_ix_i}
  \partial_\tau f(x,\tau)\,\dd x
  =\partial_\tau|I_n+\Lambda C(\tau)|^{-1/2}.
\end{equation}
Expand via \Cref{lem:det}(i) applied to $A=\Lambda C(\tau)$:
\begin{equation}\label{eq:det_expansion}
  |I_n+\Lambda C(\tau)|
  =1+\sum_{\varnothing\ne J\subseteq[n]}|(\Lambda C(\tau))_J|
  =1+\sum_{\varnothing\ne J\subseteq[n]}|C(\tau)_J|\prod_{j\in J}\lambda_j,
\end{equation}
using $|(\Lambda C(\tau))_J|=|\Lambda_J||C(\tau)_J|=\prod_{j\in J}\lambda_j\cdot|C(\tau)_J|$
since $\Lambda$ is diagonal.

\paragraph{The block structure of $|C(\tau)_J|$.}  For $J\subseteq[n]$, set
$J_1:=[n_1]\cap J$ and $J_2:=J\setminus[n_1]$.  If $J_1=\varnothing$
or $J_2=\varnothing$, then $C(\tau)_J=C_J$ is $\tau$-independent,
since the off-diagonal blocks of $C(\tau)$ (where the $\tau$ appears)
are not sampled.  Otherwise, writing $C(\tau)_J$ in block form with
blocks indexed by $J_1,J_2$,
\[
  C(\tau)_J
  =\begin{pmatrix}C_{J_1}&\tau C_{J_1J_2}\\
    \tau C_{J_2J_1}&C_{J_2}\end{pmatrix},
\]
and applying \Cref{lem:det}(ii),
\begin{equation}\label{eq:CtauJ}
  |C(\tau)_J|
  =|C_{J_1}||C_{J_2}|\prod_{i=1}^{|J_1|}(1-\tau^2\mu_{J_1,J_2}(i)),
\end{equation}
where the $\mu_{J_1,J_2}(i)\in[0,1]$ are the eigenvalues of
$C_{J_1}^{-1/2}C_{J_1J_2}C_{J_2}^{-1}C_{J_2J_1}C_{J_1}^{-1/2}$
(bounded by \Cref{lem:det}(ii) since $C_J\succ 0$).
Differentiating the product in $\tau$,
\[
  \partial_\tau|C(\tau)_J|
  =-|C_{J_1}||C_{J_2}|\sum_{i=1}^{|J_1|}2\tau\mu_{J_1,J_2}(i)
  \prod_{l\ne i}(1-\tau^2\mu_{J_1,J_2}(l))
  \leqslant 0\quad\text{on }[0,1].
\]
Define
\begin{equation}\label{eq:aJ}
  a_J(\tau):=-\partial_\tau|C(\tau)_J|\geqslant 0
  \quad(\tau\in[0,1]),
\end{equation}
with $a_J:=0$ if $J_1=\varnothing$ or $J_2=\varnothing$.  Then
\begin{equation}\label{eq:partial_det}
  \partial_\tau|I_n+\Lambda C(\tau)|
  =-\sum_{\varnothing\ne J\subseteq[n]}a_J(\tau)\prod_{j\in J}\lambda_j.
\end{equation}
Apply the chain rule
$\partial_\tau(u^{-1/2})=-\tfrac12 u^{-3/2}\,\partial_\tau u$ to
$u=|I_n+\Lambda C(\tau)|$: \eqref{eq:lt_deriv} becomes
\[
  \int_{[0,\infty)^n}e^{-\sum_i\lambda_ix_i}
    \partial_\tau f(x,\tau)\,\dd x
  =-\tfrac12|I_n+\Lambda C(\tau)|^{-3/2}
    \partial_\tau|I_n+\Lambda C(\tau)|,
\]
and substituting \eqref{eq:partial_det} (the two minus signs cancel):
\begin{equation}\label{eq:lt_of_ftau}
  \int_{[0,\infty)^n}e^{-\sum_i\lambda_ix_i}
  \partial_\tau f(x,\tau)\,\dd x
  =\sum_{\varnothing\ne J\subseteq[n]}
  \tfrac{1}{2}a_J(\tau)|I_n+\Lambda C(\tau)|^{-3/2}
  \prod_{j\in J}\lambda_j.
\end{equation}

\subsection{Identification via $h_\tau:=h_{3,C(\tau)}$}

By \Cref{lem:LaplhkC} with $k=3$ and \Cref{lem:proph}(iii) (valid
since $\alpha_i=\tfrac{3}{2}>1$),
\[
  \int e^{-\sum_i\lambda_ix_i}
    \tfrac{\partial^{|J|}}{\partial x_J}h_\tau\,\dd x
  =\Bigl(\prod_{j\in J}\lambda_j\Bigr)|I_n+\Lambda C(\tau)|^{-3/2}.
\]
Comparing with \eqref{eq:lt_of_ftau}: $\partial_\tau f(\cdot,\tau)$
and $\sum_J\tfrac12 a_J(\tau)\,\partial^{|J|}h_\tau/\partial x_J$ are
$L_1$ with the same Laplace transform, so by \Cref{thm:laplace_inj},
\begin{equation}\label{eq:key_identity}
  \partial_\tau f(x,\tau)
  =\sum_{\varnothing\ne J\subseteq[n]}
    \tfrac{1}{2}a_J(\tau)
    \tfrac{\partial^{|J|}}{\partial x_J}h_\tau(x)
  \quad\text{a.e.\ on }(0,\infty)^n.
\end{equation}

\subsection{Integration over $K$}

Integrate \eqref{eq:key_identity} over $K=\prod_j[0,s_j]$.
For $J^c:=[n]\setminus J$ and $(s_J,x_{J^c})_i=s_i$ $(i\in J)$,
$=x_i$ $(i\in J^c)$, the fundamental theorem of calculus in each
$x_j$, $j\in J$, gives
\begin{equation}\label{eq:boundary_eval}
  \int_K\tfrac{\partial^{|J|}}{\partial x_J}h_\tau(x)\,\dd x
  =\int_{\prod_{j\in J^c}[0,s_j]}h_\tau(s_J,x_{J^c})\,\dd x_{J^c}
  \geqslant 0,
\end{equation}
since the lower-endpoint contributions vanish by
\Cref{lem:proph}(ii) ($\alpha_j=\tfrac{3}{2}>1$), and
$h_\tau\geqslant 0$.  Since also $a_J\geqslant 0$,
\[
  \int_K\partial_\tau f(x,\tau)\,\dd x
  =\sum_{\varnothing\ne J\subseteq[n]}\tfrac{1}{2}a_J(\tau)
  \int_{\prod_{j\in J^c}[0,s_j]}h_\tau(s_J,x_{J^c})\,\dd x_{J^c}
  \geqslant 0,
\]
which proves \Cref{prop:key} and hence \Cref{thm:gcc2}.  \qed


%% ====================================================================
\section{\v{S}id\'ak's inequality}\label{sec:corollaries}
%% ====================================================================

\begin{corollary}[\v{S}id\'ak]\label{cor:sidak}
  For $X\sim\caln(0,C)$ and $s_i\geqslant 0$,
  \[
    \PP(|X_i|\leqslant s_i,\;i=1,\ldots,n)
    \geqslant\prod_{i=1}^n\PP(|X_i|\leqslant s_i).
  \]
\end{corollary}

\begin{proof}
  Induct on $n$.  Case $n=1$ is trivial.  For the step, take
  $K_1:=\{|x_1|\leqslant s_1\}$ and
  $L_1:=\{|x_i|\leqslant s_i,2\leqslant i\leqslant n\}$ (both closed symmetric
  convex).  \Cref{thm:gcc2} (applied with $n_1=1$, $n_2=n-1$) gives
  $\PP(X\in K_1\cap L_1)\geqslant\PP(X\in K_1)\PP(X\in L_1)$;
  apply induction to $\PP(X\in L_1)$.
\end{proof}


%% ====================================================================
\section{Worked examples}\label{sec:examples}
%% ====================================================================

We close with three exact demonstrations of the GCI on concrete
symmetric convex sets, ordered from a fully rational anchor through the
equality boundary to a genuinely correlated case in which strictness is
established in closed form.  No numerical approximation is used: every
probability is given exactly, in terms of $e$, $\pi$, and the error
function $\erf(z)=\tfrac{2}{\sqrt\pi}\int_0^z e^{-u^2}\,\dd u$.  We use
repeatedly that, for $Z\sim\caln(0,1)$,
\begin{equation}\label{eq:ex-erf}
  \PP(|Z|\leqslant a)=\Phi(a)-\Phi(-a)=2\Phi(a)-1=\erf\!\bigl(a/\sqrt2\bigr),
\end{equation}
the last equality because $\Phi(a)=\tfrac12\bigl(1+\erf(a/\sqrt2)\bigr)$.

\subsection{A fully rational anchor: nested disks}\label{sub:ex-nested}

We first record the Gaussian measure of a centred disk.  For
$G\sim\caln(0,I_2)$, passing to polar coordinates $x=(r\cos\theta,r\sin\theta)$
with Jacobian $r$,
\begin{equation}\label{eq:ex-disk}
  \mu\bigl(\{\|x\|\leqslant R\}\bigr)
  =\int_0^{2\pi}\!\!\int_0^R\frac{1}{2\pi}e^{-r^2/2}\,r\,\dd r\,\dd\theta
  =\int_0^R r\,e^{-r^2/2}\,\dd r
  =\Bigl[-e^{-r^2/2}\Bigr]_0^R
  =1-e^{-R^2/2}.
\end{equation}
(Equivalently $\|G\|^2\sim\chi^2_2$, an exponential law with mean $2$.)

\begin{example}[Concentric disks]\label{ex:nested}
  Let $G\sim\caln(0,I_2)$ and take the concentric disks
  $K=\{\|x\|\leqslant R_1\}$ and $L=\{\|x\|\leqslant R_2\}$ with
  \[
    R_1^2=2\ln 2,\qquad R_2^2=2\ln 4=4\ln 2 .
  \]
  By \eqref{eq:ex-disk},
  \[
    \mu(K)=1-e^{-\ln 2}=1-\tfrac12=\tfrac12,\qquad
    \mu(L)=1-e^{-2\ln 2}=1-\tfrac14=\tfrac34 .
  \]
  Since $R_1<R_2$ we have $K\subseteq L$, so $K\cap L=K$ and
  $\mu(K\cap L)=\tfrac12$.  The GCI \eqref{eq:gcc} reads
  \[
    \tfrac12\;\geqslant\;\tfrac12\cdot\tfrac34=\tfrac38 ,
  \]
  a strict inequality with both sides rational.
\end{example}

This is the degenerate \emph{nested} regime: whenever $K\subseteq L$ one
has $K\cap L=K$, so the GCI collapses to
$\mu(K)\geqslant\mu(K)\,\mu(L)$, i.e.\ $\mu(L)\leqslant 1$, with equality
only in the trivial case $\mu(L)=1$.  It serves purely as an exact
sanity check; the substance is in \Cref{ex:box}.

\subsection{The equality boundary: independent blocks}\label{sub:ex-indep}

\begin{example}[Slab $\times$ disk on disjoint coordinates]\label{ex:indep}
  Let $\mu=\caln(0,I_3)$ and, with $n_1=1$ and $n_2=2$, set
  \[
    K=\{x\in\R^3:|x_1|\leqslant a\},\qquad
    L=\{x\in\R^3:x_2^2+x_3^2\leqslant R^2\},
  \]
  both symmetric and convex.  Because the standard normal density on
  $\R^3$ factorises across the two coordinate blocks,
  \[
    \frac{e^{-\|x\|^2/2}}{(2\pi)^{3/2}}
    =\underbrace{\frac{e^{-x_1^2/2}}{\sqrt{2\pi}}}_{\text{block }\{1\}}
     \cdot
     \underbrace{\frac{e^{-(x_2^2+x_3^2)/2}}{2\pi}}_{\text{block }\{2,3\}},
  \]
  and $K\cap L$ is the product set
  $\{|x_1|\leqslant a\}\times\{x_2^2+x_3^2\leqslant R^2\}$, Fubini gives
  \[
    \mu(K\cap L)
    =\PP(|X_1|\leqslant a)\cdot\PP(X_2^2+X_3^2\leqslant R^2)
    =\mu(K)\,\mu(L).
  \]
  Evaluating the two factors by \eqref{eq:ex-erf} and \eqref{eq:ex-disk},
  \[
    \mu(K\cap L)=\mu(K)\,\mu(L)
    =\erf\!\Bigl(\tfrac{a}{\sqrt2}\Bigr)\,\bigl(1-e^{-R^2/2}\bigr),
  \]
  an \emph{exact equality} in \eqref{eq:gcc}.
\end{example}

This is the tight case of the GCI.  The covariance coupling the
$K$-block $\{1\}$ to the $L$-block $\{2,3\}$ is the zero block,
$C_{12}=0$, so \Cref{prop:equality} predicts equality --- as the direct
computation confirms.  Geometrically, the two sets constrain orthogonal
subspaces, and orthogonality is precisely where no correlation can be
gained.

\subsection{A correlated box: strictness in closed form}\label{sub:ex-box}

The final example is the genuinely nontrivial one.  A correlated
Gaussian box probability has no elementary closed form, but its
derivative in the correlation does, and this suffices to prove strict
inequality exactly.  The tool is the following identity.

\begin{lemma}[Plackett / Price]\label{lem:ex-plackett}
  Let $\Phi_2(h,k;\rho)$ and $\phi_2(h,k;\rho)$ denote the CDF and
  density of $\caln\!\bigl(0,\bigl(\begin{smallmatrix}1&\rho\\\rho&1
  \end{smallmatrix}\bigr)\bigr)$, $|\rho|<1$.  Then
  \begin{equation}\label{eq:ex-plackett}
    \frac{\partial}{\partial\rho}\,\Phi_2(h,k;\rho)
    =\phi_2(h,k;\rho).
  \end{equation}
\end{lemma}

\begin{proof}
  Write the density as the inverse Fourier transform of the bivariate
  normal characteristic function,
  \[
    \phi_2(h,k;\rho)
    =\frac{1}{(2\pi)^2}\iint_{\R^2}
      e^{-\mathrm i(sh+tk)}\,
      e^{-\frac12(s^2+2\rho st+t^2)}\,\dd s\,\dd t .
  \]
  Differentiating under the integral sign (justified by the Gaussian
  factor, which dominates all $s,t$-polynomials uniformly for $\rho$ in
  compact subsets of $(-1,1)$) brings down $-st$ from the exponent:
  \[
    \frac{\partial}{\partial\rho}\phi_2
    =\frac{1}{(2\pi)^2}\iint(-st)\,
      e^{-\mathrm i(sh+tk)}e^{-\frac12(\cdots)}\,\dd s\,\dd t .
  \]
  On the other hand $\partial_h$ and $\partial_k$ bring down $-\mathrm i s$
  and $-\mathrm i t$ respectively, so
  $\partial_h\partial_k\phi_2$ produces the factor
  $(-\mathrm i s)(-\mathrm i t)=-st$ under the same integral.  Hence
  \begin{equation}\label{eq:ex-heat}
    \frac{\partial\phi_2}{\partial\rho}
    =\frac{\partial^2\phi_2}{\partial h\,\partial k}.
  \end{equation}
  Now integrate \eqref{eq:ex-heat} over the quadrant
  $(-\infty,h]\times(-\infty,k]$.  Since $\phi_2$ and its first partials
  decay to $0$ at $-\infty$, two applications of the fundamental theorem
  of calculus give
  \[
    \frac{\partial}{\partial\rho}\Phi_2(h,k;\rho)
    =\int_{-\infty}^{h}\!\!\int_{-\infty}^{k}
      \frac{\partial\phi_2}{\partial\rho}\,\dd v\,\dd u
    =\int_{-\infty}^{h}\!\!\int_{-\infty}^{k}
      \frac{\partial^2\phi_2}{\partial u\,\partial v}\,\dd v\,\dd u
    =\phi_2(h,k;\rho).\qedhere
  \]
\end{proof}

\begin{example}[Two correlated slabs; \v{S}id\'ak with $n_1=n_2=1$]\label{ex:box}
  Let $X\sim\caln(0,C_\rho)$ on $\R^2$ with
  $C_\rho=\bigl(\begin{smallmatrix}1&\rho\\\rho&1\end{smallmatrix}\bigr)$,
  $\rho\in(-1,1)$, and take the symmetric slabs
  \[
    K=\{|x_1|\leqslant a\},\qquad L=\{|x_2|\leqslant a\},\qquad a>0,
  \]
  so that $K\cap L=[-a,a]^2$.  Put
  $P(\rho):=\mu(K\cap L)=\PP\bigl(|X_1|\leqslant a,\,|X_2|\leqslant a\bigr)$.
  Then
  \begin{equation}\label{eq:ex-gap}
    \mu(K\cap L)-\mu(K)\,\mu(L)
    =\int_{0}^{|\rho|}
      \frac{e^{-a^2/(1+s)}-e^{-a^2/(1-s)}}{\pi\sqrt{1-s^2}}\,\dd s
    \;>\;0
    \qquad(\rho\neq 0),
  \end{equation}
  so the GCI \eqref{eq:gcc} holds \emph{strictly} for every $\rho\neq0$,
  with equality iff $\rho=0$.
\end{example}

\begin{proof}
  By inclusion--exclusion for the rectangle $[-a,a]^2$, writing
  $\Phi_2(h,k)=\PP(X_1\leqslant h,\,X_2\leqslant k)$,
  \[
    P(\rho)=\Phi_2(a,a)-\Phi_2(-a,a)-\Phi_2(a,-a)+\Phi_2(-a,-a).
  \]
  Differentiating in $\rho$ and applying \Cref{lem:ex-plackett} to each
  term,
  \begin{equation}\label{eq:ex-Pprime0}
    P'(\rho)=\phi_2(a,a;\rho)-\phi_2(-a,a;\rho)-\phi_2(a,-a;\rho)
      +\phi_2(-a,-a;\rho).
  \end{equation}
  The exponent of $\phi_2$ is
  $Q(h,k)=\dfrac{h^2-2\rho hk+k^2}{2(1-\rho^2)}$.  At the four corners,
  \[
    Q(\pm a,\pm a)=\frac{2a^2-2\rho a^2}{2(1-\rho^2)}
      =\frac{a^2(1-\rho)}{(1-\rho)(1+\rho)}=\frac{a^2}{1+\rho},
    \qquad
    Q(\pm a,\mp a)=\frac{2a^2+2\rho a^2}{2(1-\rho^2)}=\frac{a^2}{1-\rho},
  \]
  (each value occurring twice, by $Q(h,k)=Q(-h,-k)$).  Hence the
  like-sign corners share the value
  $\frac{1}{2\pi\sqrt{1-\rho^2}}e^{-a^2/(1+\rho)}$ and the opposite-sign
  corners the value $\frac{1}{2\pi\sqrt{1-\rho^2}}e^{-a^2/(1-\rho)}$, and
  \eqref{eq:ex-Pprime0} collapses to
  \begin{equation}\label{eq:ex-Pprime}
    P'(\rho)=\frac{1}{\pi\sqrt{1-\rho^2}}
      \Bigl(e^{-a^2/(1+\rho)}-e^{-a^2/(1-\rho)}\Bigr).
  \end{equation}
  For $\rho\in(0,1)$ we have $0<1-\rho<1+\rho$, so
  $\frac{a^2}{1+\rho}<\frac{a^2}{1-\rho}$ and therefore
  $e^{-a^2/(1+\rho)}>e^{-a^2/(1-\rho)}$, giving $P'(\rho)>0$.
  Moreover the reflection $(x_1,x_2)\mapsto(x_1,-x_2)$ preserves both
  Lebesgue measure and the box $[-a,a]^2$ while sending $C_\rho$ to
  $C_{-\rho}$; hence $P(\rho)=P(-\rho)$, and $P$ is even.  Thus $P$ is
  strictly increasing on $[0,1)$ and strictly decreasing on $(-1,0]$,
  with its unique minimum at $\rho=0$.

  At $\rho=0$ the coordinates are independent, so by \eqref{eq:ex-erf}
  \[
    P(0)=\PP(|X_1|\leqslant a)\,\PP(|X_2|\leqslant a)
        =\erf\!\Bigl(\tfrac{a}{\sqrt2}\Bigr)^{2}=\mu(K)\,\mu(L).
  \]
  Integrating \eqref{eq:ex-Pprime} from $0$ to $\rho>0$ (and using
  evenness for $\rho<0$) gives
  $P(\rho)-P(0)=\int_0^{|\rho|}P'(s)\,\dd s$, which is
  \eqref{eq:ex-gap}; the integrand is strictly positive on $(0,|\rho|]$,
  so the gap is strictly positive.
\end{proof}

\begin{remark}[Exact size of the effect as $a\to 0$]\label{rem:ex-smalla}
  Expanding the integrand of \eqref{eq:ex-gap} for small $a$ via
  $e^{-a^2/(1\pm s)}=1-\tfrac{a^2}{1\pm s}+O(a^4)$ gives
  \[
    e^{-a^2/(1+s)}-e^{-a^2/(1-s)}
    =a^2\Bigl(\tfrac{1}{1-s}-\tfrac{1}{1+s}\Bigr)+O(a^4)
    =\frac{2a^2 s}{1-s^2}+O(a^4),
  \]
  so, with the substitution $u=1-s^2$ in
  $\int_0^{\rho}s(1-s^2)^{-3/2}\,\dd s=\bigl[(1-s^2)^{-1/2}\bigr]_0^{\rho}
  =(1-\rho^2)^{-1/2}-1$,
  \[
    \mu(K\cap L)-\mu(K)\,\mu(L)
    =\frac{2a^2}{\pi}\Bigl((1-\rho^2)^{-1/2}-1\Bigr)+O(a^4).
  \]
  Since $\erf(a/\sqrt2)=\sqrt{2/\pi}\,a+O(a^3)$ gives
  $\mu(K)\,\mu(L)=\tfrac{2a^2}{\pi}+O(a^4)$, the relative discrepancy
  tends to an exact constant,
  \[
    \frac{\mu(K\cap L)}{\mu(K)\,\mu(L)}
    \;\xrightarrow[a\to0]{}\;\frac{1}{\sqrt{1-\rho^2}}\;>\;1
    \qquad(\rho\neq0),
  \]
  quantifying exactly how far the inequality is from tight for small
  boxes.
\end{remark}

\begin{remark}[Geometric reading and the general \v{S}id\'ak inequality]
  \label{rem:ex-geom}
  Under the rotation-invariant law $\caln(0,I_2)$, two slabs of
  half-width $a$ whose unit normals $u_1,u_2$ meet at angle $\theta$
  define the forms $X_i=\langle u_i,G\rangle\sim\caln(0,1)$ with
  $\Cov(X_1,X_2)=\langle u_1,u_2\rangle=\cos\theta$.  Thus \Cref{ex:box}
  with $\rho=\cos\theta$ says that the rhombic cross-strip $K\cap L$
  always carries at least the product measure, strictly so unless
  $\theta=\tfrac\pi2$; perpendicular slabs are independent and recover
  the equality of \Cref{ex:indep}.  \Cref{ex:box} is moreover a
  self-contained, elementary proof of \Cref{cor:sidak} in the case
  $n=2$, independent of the machinery of \Cref{sec:proof}; the general
  $n$ case is the content of \Cref{cor:sidak} (via \Cref{thm:gcc2}), and
  was first obtained by \v{S}id\'ak (1967) through a related
  differentiation in the correlations.
\end{remark}



\subsection{Statistical applications}\label{sub:ex-stats}

The inequality was conjectured and first studied because of the
multiple-inference problem below, where it furnishes guarantees that
hold \emph{uniformly over the unknown dependence structure}.  We give
three worked applications: dependence-free simultaneous inference, an
exact measure of its conservativeness, and the extension to estimated
scale (the multivariate $t$).  Here $\Phi$ and $\Phi^{-1}$ are the
standard normal CDF and quantile.

\begin{example}[\v{S}id\'ak simultaneous intervals and familywise error]
  \label{ex:sidak-ci}
  Let $\hat\theta\sim\caln(\theta,\Sigma)$ in $\R^m$ with known variances
  $\sigma_i^2:=\Sigma_{ii}>0$ and \emph{arbitrary} correlation matrix
  $R$.  Standardizing, $Z:=\bigl((\hat\theta_i-\theta_i)/\sigma_i\bigr)_i
  \sim\caln(0,R)$ with each $Z_i\sim\caln(0,1)$.  For $c>0$ form the
  rectangle $\mathcal R_c:=\prod_{i=1}^m[\hat\theta_i-c\sigma_i,\,
  \hat\theta_i+c\sigma_i]$.  Then
  \begin{equation}\label{eq:ex-sidak-cov}
    \PP(\theta\in\mathcal R_c)
    =\PP\Bigl(\textstyle\bigcap_{i=1}^m\{|Z_i|\leqslant c\}\Bigr)
    \;\geqslant\;\prod_{i=1}^m\PP(|Z_i|\leqslant c)
    =\bigl(2\Phi(c)-1\bigr)^m ,
  \end{equation}
  for every $R$.  Hence the \v{S}id\'ak choice
  \[
    c=c_S:=\Phi^{-1}\!\Bigl(\tfrac{1+(1-\alpha)^{1/m}}{2}\Bigr)
  \]
  yields simultaneous coverage $\geqslant 1-\alpha$; dually, testing
  $H_{0,i}:\theta_i=0$ by rejecting when $|Z_i|>c_S$ controls the
  familywise error rate (FWER) at level $\leqslant\alpha$ irrespective of
  $R$.
\end{example}

\begin{proof}
  The sets $K_i=\{z:|z_i|\leqslant c\}$ are symmetric slabs, and
  $|\hat\theta_i-\theta_i|\leqslant c\sigma_i\iff|Z_i|\leqslant c$, so the
  coverage event is $\bigcap_i\{|Z_i|\leqslant c\}$ and
  \eqref{eq:ex-sidak-cov} is exactly \Cref{cor:sidak}.  At $c=c_S$ one
  has $2\Phi(c_S)-1=(1-\alpha)^{1/m}$, so the lower bound equals
  $1-\alpha$.  For the testing statement, under $\theta=0$,
  \[
    \mathrm{FWER}
    =\PP\Bigl(\textstyle\bigcup_i\{|Z_i|>c\}\Bigr)
    =1-\PP\Bigl(\textstyle\bigcap_i\{|Z_i|\leqslant c\}\Bigr)
    \leqslant 1-\bigl(2\Phi(c)-1\bigr)^m ,
  \]
  which is $\leqslant\alpha$ at $c=c_S$.
\end{proof}

\begin{remark}[Uniformly sharper than Bonferroni]\label{rem:ex-bonf}
  The Bonferroni rule uses $c_B:=\Phi^{-1}\bigl(1-\tfrac{\alpha}{2m}\bigr)$,
  so $2\Phi(c_B)-1=1-\tfrac{\alpha}{m}$.  By Bernoulli's inequality
  $(1-\alpha)^{1/m}\leqslant 1-\tfrac{\alpha}{m}$ (valid since
  $0<\tfrac1m\leqslant 1$, strict for $m\geqslant 2$),
  \[
    2\Phi(c_S)-1=(1-\alpha)^{1/m}\leqslant 1-\tfrac{\alpha}{m}=2\Phi(c_B)-1,
  \]
  whence $c_S\leqslant c_B$.  The \v{S}id\'ak rectangle is therefore
  contained in the Bonferroni rectangle and gives strictly shorter
  intervals (equivalently, a uniformly more powerful test) at the same
  guaranteed level for $m\geqslant 2$.
\end{remark}

\begin{example}[Exact conservativeness for two correlated tests]
  \label{ex:fwer-exact}
  Take $m=2$ with $\Cov(Z_1,Z_2)=\rho$ (so the correlation is $\rho$,
  the marginals being standard).  The true simultaneous coverage at
  threshold $c$ is $P(\rho)$ from \Cref{ex:box} with $a=c$, while the
  dependence-free guarantee \eqref{eq:ex-sidak-cov} is
  $P(0)=(2\Phi(c)-1)^2$.  By \eqref{eq:ex-gap} the excess coverage ---
  equivalently, the amount by which the actual FWER falls below the
  nominal level --- is exactly
  \[
    \PP(|Z_1|\leqslant c,\,|Z_2|\leqslant c)-(2\Phi(c)-1)^2
    =\int_0^{|\rho|}\frac{e^{-c^2/(1+s)}-e^{-c^2/(1-s)}}
      {\pi\sqrt{1-s^2}}\,\dd s\;\geqslant\;0 .
  \]
  The procedure is thus conservative under any dependence, with slack
  increasing in $|\rho|$.  Two regimes are explicit: as $c\to0$ (a very
  stringent level) the ratio of true to guaranteed coverage tends to
  $(1-\rho^2)^{-1/2}$ by \Cref{rem:ex-smalla}; and as $\rho\to1^-$
  (perfectly dependent statistics) the two events coincide,
  $P(\rho)\to 2\Phi(c)-1$, so the excess attains its maximum
  \[
    \bigl(2\Phi(c)-1\bigr)-\bigl(2\Phi(c)-1\bigr)^2
    =2\bigl(2\Phi(c)-1\bigr)\bigl(1-\Phi(c)\bigr).
  \]
  This makes precise the familiar fact that \v{S}id\'ak (and a fortiori
  Bonferroni) corrections are most conservative when the test statistics
  are strongly correlated.
\end{example}

\begin{example}[Unknown variance: \v{S}id\'ak for the multivariate $t$]
  \label{ex:mvt}
  When the scale is estimated, the standardized statistics follow a
  multivariate $t$.  Model $X=(Z_1/S,\ldots,Z_m/S)$ with $Z\sim\caln(0,R)$
  ($R$ a correlation matrix) and $S>0$ independent of $Z$; if $\nu
  S^2\sim\chi^2_\nu$ then $X$ has the central $t_\nu(R)$ law and each
  $X_i\sim t_\nu$.  Writing $F_\nu$ for the $t_\nu$ CDF, for any
  thresholds $c_1,\ldots,c_m>0$,
  \begin{equation}\label{eq:ex-mvt}
    \PP\Bigl(\textstyle\bigcap_{i=1}^m\{|X_i|\leqslant c_i\}\Bigr)
    \;\geqslant\;\prod_{i=1}^m\PP(|X_i|\leqslant c_i)
    =\prod_{i=1}^m\bigl(2F_\nu(c_i)-1\bigr).
  \end{equation}
  In particular the \v{S}id\'ak threshold built from $t_\nu$-quantiles
  controls simultaneous coverage and FWER for correlated $t$-statistics
  under arbitrary $R$.
\end{example}

\begin{proof}
  Condition on $S$.  For fixed $s>0$, $\{|X_i|\leqslant c_i\}=\{|Z_i|
  \leqslant sc_i\}$, so the Gaussian inequality \Cref{cor:sidak},
  applied to the slabs $\{|z_i|\leqslant sc_i\}$ and using
  $Z_i\sim\caln(0,1)$, gives
  \[
    \PP\Bigl(\textstyle\bigcap_i\{|X_i|\leqslant c_i\}\,\Big|\,S=s\Bigr)
    \;\geqslant\;\prod_{i=1}^m\bigl(2\Phi(sc_i)-1\bigr)
    =:\prod_{i=1}^m g_i(s).
  \]
  Taking expectations over $S$,
  \[
    \PP\Bigl(\textstyle\bigcap_i\{|X_i|\leqslant c_i\}\Bigr)
    =\EE\Bigl[\PP\bigl(\textstyle\bigcap_i\{|Z_i|\leqslant Sc_i\}\mid S\bigr)\Bigr]
    \;\geqslant\;\EE\Bigl[\prod_{i=1}^m g_i(S)\Bigr].
  \]
  Each $g_i(s)=2\Phi(sc_i)-1$ is nonnegative and nondecreasing on
  $(0,\infty)$, so the $g_i$ are simultaneously nondecreasing functions
  of the single variable $S$.  A product of nonnegative nondecreasing
  functions is again nondecreasing, so Chebyshev's association
  inequality $\EE[fg]\geqslant\EE[f]\,\EE[g]$ (for nondecreasing $f,g$ of
  one variable), applied iteratively, yields
  $\EE[\prod_i g_i(S)]\geqslant\prod_i\EE[g_i(S)]$.  Finally
  $\EE[g_i(S)]=\EE\bigl[2\Phi(Sc_i)-1\bigr]
  =\PP(|Z_i|\leqslant Sc_i)=\PP(|X_i|\leqslant c_i)=2F_\nu(c_i)-1$.
  Combining the three displays gives \eqref{eq:ex-mvt}.
\end{proof}

\begin{remark}
  The three examples isolate distinct virtues of the GCI for inference:
  \Cref{ex:sidak-ci} gives coverage and error control that are
  \emph{free of the dependence} $R$; \Cref{ex:fwer-exact} measures the
  resulting conservativeness exactly; and \Cref{ex:mvt} shows the same
  guarantees survive estimation of the scale, the case of practical
  relevance for $t$-statistics.  Notably the last step combines the GCI
  with the classical Chebyshev correlation inequality.
\end{remark}


%% ====================================================================
\section{Notes}\label{sec:notes}
%% ====================================================================

\subsection{Historical remarks}\label{note:history}

The conjecture originates in statistical work on joint coverage of
rectangular confidence regions for multivariate normal means:
Dunnett--Sobel (1955), Dunn (1958).  Das Gupta, Eaton, Olkin,
Perlman, Savage, and Sobel (1972) stated it in geometric form.
Partial results before Royen:
\begin{itemize}[nosep]
  \item Khatri (1967) and \v{S}id\'ak (1967), independently: one set
    is a symmetric strip (= intersection of symmetric slabs).  This
    is our \Cref{cor:sidak}.
  \item Pitt (1977): $d=2$.
  \item Schechtman--Schlumprecht--Zinn (1998): both sets are
    centered ellipsoids.
  \item Harg\'e (1999): one set is a centered ellipsoid.
\end{itemize}
Royen (2014) proved the full conjecture and, more generally, the
analogous inequality for multivariate gamma vectors
$\Gamma_n(\alpha,R)$ with $\alpha>0$ half-integer; the Gaussian case
is $\alpha=\tfrac{1}{2}$.  Lata{\l}a--Matlak (2015) specialised and
reorganised the argument for the Gaussian case, which is the version
followed here.

\subsection{Extension to singular covariance}\label{note:singular}

Let $C\succcurlyeq 0$ with $\rank C=r\leqslant n$.  By the spectral theorem,
$C=O\diag(\lambda_1,\ldots,\lambda_r,0,\ldots,0)O^\top$ with $O$
orthogonal and $\lambda_i>0$.  Set $B\in M_{n\times r}$ to have the
$i$-th column equal to $\sqrt{\lambda_i}\,O_{\bullet i}$; then
$BB^\top=C$ and $B$ has full column rank.  If $G\sim\caln(0,I_r)$,
then $BG\sim\caln(0,BB^\top)=\caln(0,C)$.

For closed symmetric convex $K\subseteq\R^n$:
\begin{itemize}[nosep]
  \item $B^{-1}(K)=\{g\in\R^r:Bg\in K\}$ is closed by continuity of
    $g\mapsto Bg$;
  \item convex, since $Bg_1,Bg_2\in K$ implies
    $B(\alpha g_1+(1-\alpha)g_2)=\alpha Bg_1+(1-\alpha)Bg_2\in K$ by
    convexity of $K$;
  \item symmetric, since $g\in B^{-1}(K)\iff Bg\in K\iff -Bg\in K\iff
    -g\in B^{-1}(K)$.
\end{itemize}
Hence, writing $\mu=\caln(0,C)$ and $\nu=\caln(0,I_r)$,
\[
  \mu(K\cap L)
  =\PP(BG\in K\cap L)
  =\nu(B^{-1}(K)\cap B^{-1}(L))
  \geqslant\nu(B^{-1}(K))\nu(B^{-1}(L))
  =\mu(K)\mu(L),
\]
where the inequality is \Cref{thm:gcc_geom} in dimension $r$
(non-singular case, already established).

\subsection{The choice $k=3$}\label{note:why_k3}

Two constraints pin down the shape parameter.
\begin{enumerate}
  \item Laplace-side matching.  Differentiating
    $|I+\Lambda C(\tau)|^{-1/2}$ in $\tau$ yields a linear combination
    of $|I+\Lambda C(\tau)|^{-3/2}\prod_{j\in J}\lambda_j$
    (\eqref{eq:lt_of_ftau}).  To rewrite this as the Laplace transform
    of a density, we need an auxiliary $h_\tau$ whose Laplace
    transform is $|I+\Lambda C(\tau)|^{-3/2}$, i.e.,
    $h_\tau=h_{k,C(\tau)}$ with $k=3$ in \Cref{lem:LaplhkC}.

  \item Integration by parts.  \Cref{lem:proph}(ii)--(iii) require
    $\alpha_i=k/2>1$ both for $h_\alpha(x)\to 0$ as $x_i\to 0+$ and for
    validity of the derivative identity
    $\partial_xg_\alpha=g_{\alpha-1}-g_\alpha$.  $k=3$ gives
    $\alpha_i=\tfrac{3}{2}>1$.
\end{enumerate}

The method applies whenever the Laplace transform of the
underlying squared vector has the form $|I+\Lambda C|^{-\beta}$ with
$\beta+1>1$, i.e., $\beta>0$; in the Gaussian case $\beta=\tfrac{1}{2}$.
This is the range covered by Royen's multivariate gamma result.

\subsection{Strict positivity of $a_J$}\label{note:phase}

Under $C\succ 0$, the matrix
$C_{J_1}^{-1/2}C_{J_1J_2}C_{J_2}^{-1}C_{J_2J_1}C_{J_1}^{-1/2}$ is
the matrix of squared canonical correlations between the coordinate
blocks $J_1$ and $J_2$.  Its eigenvalues $\mu_{J_1,J_2}(i)$ lie in
$[0,1)$: evaluating \eqref{eq:CtauJ} at $\tau=1$ gives
$|C_J|=|C_{J_1}||C_{J_2}|\prod_i\bigl(1-\mu_{J_1,J_2}(i)\bigr)$, and since
$C\succ 0$ forces $|C_J|,|C_{J_1}|,|C_{J_2}|>0$, every factor
$1-\mu_{J_1,J_2}(i)$ is strictly positive, i.e.\ $\mu_{J_1,J_2}(i)<1$.
Consequently, for any $J$ with
$J_1,J_2\ne\varnothing$ and $C_{J_1J_2}\ne 0$, some $\mu_{J_1,J_2}(i)
>0$, so $a_J(\tau)>0$ on $(0,1]$.  At $\tau=0$, every $a_J(0)=0$
because each $|C(\tau)_J|$ is even in $\tau$ (see \eqref{eq:CtauJ}).

\subsection{Equality in the GCI}\label{note:equality}

\begin{proposition}\label{prop:equality}
  Let $X\sim\caln(0,C)$, $C\succ 0$, and let $K,L\subseteq\R^n$ be as
  in \Cref{thm:gcc2}.  Then equality
  $\PP(X\in K\cap L)=\PP(X\in K)\PP(X\in L)$ holds iff $C_{12}=0$.
\end{proposition}

\begin{proof}
  If $C_{12}=0$, then $C(\tau)=\diag(C_{11},C_{22})$ for all $\tau$,
  so the two coordinate blocks of $X$ are independent and both sides
  of \eqref{eq:gcc2} are equal.

  Conversely, suppose equality.  From the fundamental theorem of
  calculus and \eqref{eq:key_identity}:
  \[
    \PP(X\in K\cap L)-\PP(X\in K)\PP(X\in L)
    =\int_0^1\!\!\sum_{J}\tfrac{1}{2}a_J(\tau)
    \int_{\prod_{J^c}[0,s_j]}h_\tau(s_J,x_{J^c})\,\dd x_{J^c}\,\dd\tau.
  \]
  Each $a_J$ is polynomial in $\tau$ (since $|C(\tau)_J|$ is
  polynomial in the entries, which are affine in $\tau$), hence
  continuous on $[0,1]$ with $a_J\geqslant 0$.  The inner integral is $>0$
  for every $\tau\in[0,1]$: $g_{3/2}(\cdot,y)>0$ on $(0,\infty)$ for
  each $y\geqslant 0$, because every term in the series of
  \Cref{def:ncgamma} is nonnegative and the $k=0$ term equals
  $x^{1/2}e^{-x-y}/\Gamma(3/2)>0$; taking expectation over $Y$
  preserves strict positivity, so $h_\tau>0$ on $(0,\infty)^n$ and the
  integral over any product $\prod_{J^c}[0,s_j]$ of positive lengths
  is $>0$.

  Equality thus forces $a_J(\tau)\equiv 0$ on $[0,1]$ (by continuity)
  for every $J$ with $J_1,J_2\ne\varnothing$.  By \Cref{note:phase},
  this means $C_{J_1J_2}=0$ for every such $J$.  Taking $J=\{i,j\}$
  with $i\leqslant n_1<j\leqslant n$ yields $C_{ij}=0$, so $C_{12}=0$.
\end{proof}

\subsection{Verification of \eqref{eq:boundary_eval} for $n=2$}
  \label{note:boundary}

For $J=\{1\}$:
\[
  \int_0^{s_1}\!\!\int_0^{s_2}\partial_{x_1}h_\tau\,\dd x_1\,\dd x_2
  =\int_0^{s_2}\bigl(h_\tau(s_1,x_2)-h_\tau(0+,x_2)\bigr)\dd x_2
  =\int_0^{s_2}h_\tau(s_1,x_2)\,\dd x_2.
\]
For $J=\{1,2\}$:
$\int_0^{s_1}\!\int_0^{s_2}\partial^2_{x_1x_2}h_\tau=h_\tau(s_1,s_2)
-h_\tau(s_1,0+)-h_\tau(0+,s_2)+h_\tau(0+,0+)=h_\tau(s_1,s_2)$.
Both use $h_\tau(\cdot)\to 0$ at $0+$ from \Cref{lem:proph}(ii).

\subsection{Necessity of symmetry}\label{note:symmetry}

Without $K=-K$, $L=-L$ the inequality fails.  Under $\mu=\caln(0,I_2)$
take
\[
  K:=\{x:x_1\geqslant 1\},
  \qquad
  L:=\{x:x_1\leqslant -1\}.
\]
Both are closed convex, neither is symmetric.  Then
$\mu(K)=\mu(L)=1-\Phi(1)\approx 0.1587$, so
$\mu(K)\mu(L)\approx 0.0252$, while $K\cap L=\varnothing$ gives
$\mu(K\cap L)=0<\mu(K)\mu(L)$.

\subsection{\v{S}id\'ak correction for multiple testing}\label{note:fwer}

The use of \Cref{cor:sidak} to obtain dependence-free simultaneous
confidence intervals and familywise error control, its comparison with
the Bonferroni correction, the exact conservativeness in the bivariate
case, and the extension to estimated scale (the multivariate $t$) are
developed in \Cref{sub:ex-stats} (\Cref{ex:sidak-ci,ex:fwer-exact,ex:mvt}).

\subsection{Related correlation inequalities}\label{note:correlations}

\begin{center}
\renewcommand{\arraystretch}{1.3}
\begin{tabular}{llll}
\hline
 & \textbf{Measure} & \textbf{Events/functions} & \textbf{Conclusion} \\
\hline
GCI & Centered Gaussian on $\R^n$ & Symmetric convex $K,L$
 & $\mu(K\cap L)\geqslant\mu(K)\mu(L)$ \\
FKG (1971) & Log-concave on a lattice & Monotone $f,g$
 & $\EE[fg]\geqslant\EE f\cdot\EE g$ \\
Harris (1960) & Product on $\{0,1\}^n$ & Monotone events $A,B$
 & $\PP(A\cap B)\geqslant\PP(A)\PP(B)$ \\
\hline
\end{tabular}
\end{center}

FKG contains Harris as a special case.  The GCI is formally
unrelated to FKG: symmetric convex sets in $\R^n$ are not monotone
for any coordinate order, and monotone events on a lattice are not
symmetric.

\subsection{Open questions}\label{note:open}

\begin{enumerate}
  \item Quantitative GCI.  No explicit nontrivial lower bound on the
    excess $\mu(K\cap L)-\mu(K)\mu(L)$ is known in the generality of
    \Cref{thm:gcc_geom}.  Candidates depending on Gaussian widths,
    eccentricity, and the spectral gap of $C$ have been proposed but
    not established.
  \item Non-Gaussian extensions.  Counterexamples rule out general
    log-concave measures.  The exact class of product-type
    inequalities valid for symmetric convex sets — beyond the Gaussian
    and multivariate-gamma families of Royen — is open.
  \item Functional forms.  Whether
    $\EE[f(X)g(X)]\geqslant\EE f(X)\,\EE g(X)$ holds for centered
    Gaussian $X$ and even quasi-concave $f,g\geqslant 0$ is open;
    partial results are in Lata\l{}a--Matlak (2015, Remark~3).
\end{enumerate}


%% ====================================================================
\section*{References}
%% ====================================================================

\begin{enumerate}[nosep,label={[\arabic*]}]
  \item T.\ Royen, ``A simple proof of the Gaussian correlation
    conjecture extended to multivariate gamma distributions,''
    \textit{Far East J.\ Theor.\ Stat.}\ \textbf{48} (2014), 139--145.
    \texttt{arXiv:1408.1028}.
  \item R.\ Lata{\l}a and D.\ Matlak, ``Royen's proof of the Gaussian
    correlation inequality,'' in \textit{Geometric Aspects of Functional
    Analysis}, Lecture Notes in Math.\ \textbf{2169}, Springer, 2017,
    pp.\ 265--275.  \texttt{arXiv:1512.08776}.
  \item S.\ Das Gupta, M.\ L.\ Eaton, I.\ Olkin, M.\ Perlman,
    L.\ J.\ Savage, and M.\ Sobel, ``Inequalities on the probability
    content of convex regions for elliptically contoured
    distributions,'' \textit{Proc.\ Sixth Berkeley Symp.}, 1972.
  \item L.\ D.\ Pitt, ``A Gaussian correlation inequality for symmetric
    convex sets,'' \textit{Ann.\ Probab.}\ \textbf{5} (1977), 470--474.
  \item C.\ G.\ Khatri, ``On certain inequalities for normal
    distributions,'' \textit{Ann.\ Math.\ Statist.}\ \textbf{38}
    (1967), 1853--1867.
  \item Z.\ \v{S}id\'ak, ``Rectangular confidence regions for the means
    of multivariate normal distributions,'' \textit{J.\ Amer.\ Statist.\
    Assoc.}\ \textbf{62} (1967), 626--633.
  \item G.\ Schechtman, T.\ Schlumprecht, and J.\ Zinn, ``On the
    Gaussian measure of the intersection,'' \textit{Ann.\ Probab.}\
    \textbf{26} (1998), 346--357.
  \item G.\ Harg\'e, ``A particular case of correlation inequality for
    the Gaussian measure,'' \textit{Ann.\ Probab.}\ \textbf{27} (1999),
    1939--1951.
  \item C.\ Borell, ``A Gaussian correlation inequality for certain
    bodies in $\R^n$,'' \textit{Math.\ Ann.}\ \textbf{256} (1981),
    569--573.
\end{enumerate}

\end{document}
